% id: PHY-0036
% title: เวลาเมื่ออัตราเร็วขาลงเป็นครึ่งหนึ่งของอัตราเร็วต้น
% camp: อื่น ๆ
% topic: การเคลื่อนที่แนวตรง
% year: 2553
% tags: IJSO, รอบที่ 1, การโยนในแนวดิ่ง, กราฟความเร็ว–เวลา
% origin: IJSO รอบที่ 1 2553 · ข้อ 36 ในภาพที่ผู้ใช้ส่ง
% review: ถอดจากภาพแล้ว · รอตรวจทาน
% diagram: PHY-0036.svg
% scan: sources/PHY-0036-original.png
\documentclass[12pt]{article}
\usepackage{amsmath,amssymb,fontspec,polyglossia,tikz}
\setdefaultlanguage{thai}
\setmainfont{Noto Serif Thai}
\usetikzlibrary{arrows.meta}
\begin{document}
\textbf{ข้อ 36.} (IJSO รอบที่ 1 2553)
อนุภาคหนึ่งถูกโยนขึ้นในแนวดิ่งภายใต้สนามโน้มถ่วง $g$ กราฟความเร็ว–เวลาเป็นดังรูป (ทิศขึ้นแทนด้วยเครื่องหมายบวก) ตอนขาลง เมื่ออนุภาคมีอัตราเร็วครึ่งหนึ่งของอัตราเร็วต้น เวลาผ่านไปเท่าใดนับจากเริ่มโยน

\begin{tikzpicture}[x=1.24533cm,y=0.52456cm,>=Stealth]
\draw[->,line width=.7pt] (-0.18,0) -- (5,0);
\draw[->,line width=.7pt] (0,-2.8) -- (0,4.8);
\draw[line width=.7pt] (2,-0.1271) -- (2,0.1271);
\node[anchor=center] at (2,-0.5084) {$T$};
\node[anchor=east] at (-0.0857,-0.4067) {0};
\node[anchor=west] at (5.0857,-0.0508) {เวลา};
\node[anchor=west] at (-0.0857,5.2067) {ความเร็ว};
\draw[line width=.7pt] (0,3) -- (2,0) -- (3.6,-2.4);
\end{tikzpicture}


1. $\dfrac{T}{2}$

2. $T\left(1+\dfrac{1}{\sqrt{2}}\right)$

3. $\dfrac{3T}{2}$

4. $T\left(2-\dfrac{1}{\sqrt{2}}\right)$
\end{document}
